\subsection{Tips til at arrangere}
\subsubsection{Inspiration}
Prøv at søge efter live-udgaver af nummeret. Nogle kunstnere gør meget for at arrangere deres numre til live-optrædener. Beyoncé og Bruno Mars er GOATede på den front.

Nogle gange har de andre revybands også gode idéer, der er veljusterede til revy-formatet, dets muligheder og begrænsninger.

\subsubsection{Vision}
De fleste revyster har ikke så meget (musikalsk) vision. Typisk har instruktør/sangskriver ikke andre idéer om hvordan \textit{revysangen} skal lyde, end at melodien bliver sunget. Det betyder ikke at man skal holde sig indenfor linierne, tværtimod; i min erfaring bliver de fleste begejstrede for kunstnerisk input og ideer. Hvis du planlægger at genrebende, bør du kommunikere det til alle der er kreativt involverede med nummeret.

\subsubsection{Gør det nemt}
Man skal lære meget materiale til en revy. Særligt tiden med samlet band er knap. Desto hurtigere man kan lære arrangementet grundlæggende, desto mere tid har man til at gøre det fedt.

\subsubsection{Kommuniker med koreografen}
Hvis der er koreografi til sangen, så vær sikker på at arrangøren og koreografen er \textit{helt} enige om formen, såvel som tempoet. Send gerne en optagelse af hvordan i spiller nummeret til dem, så de kan tjekke at deres koreografi passer til. Hvis du som arrangør planlægger at ændre formen i forhold til referencen, så kommunikér det tidligt.

\subsubsection{Blæser-arrangementer}
Unisont er godt! En blæsergruppe har naturligt en meget rigere klang end elguitar/keyboards, så der skal mindre harmonisk tæthed til, for at det "lyder af noget".

Rytmisk musik er ikke Bach. Koralharmonisering som du har lært det i gymnasiet er ikke altid fedt. Prøv i stedet at arbejde med riffs og rytmiske idéer hvor blæsersektionen har interessante rytmer og modsvar til melodien.

Nogle gange er flydetoner fint.

\subsubsection{Barytonsaxen}
SaTyR-revyen har en barytonsax (den burde ligge i brandskabet i dyrestalden). Brug den! Hvis dit band har en træblæser, kan de snildt spille på den. Barytonens \textbf{BÅT} fungerer virkelig godt i en dum-tak konstruktion med de lysere blæsere (Eksempel: c-stykket i Kom lad os gå med Det Brune Punktum). Et simpelt baa-dap ind i hver periode fungerer som mini-basfill (Eksempel: Side To Side med Ariana Grande). The Rocket med Blast inkorporerer begge idéer.

\subsubsection{Lav en demo}
Det er ofte meget nemmere at lære et arrangement, hvis man ved hvordan det skal lyde, fremfor at få det forklaret. Her kan det være nyttigt at lave en demo. Med lidt øvelse kan det gøres relativt nemt i en DAW. Brug referencen som udgangspunkt og byg på. Brug evt. Moises eller lignende til at separære tracks hvis nødvendigt. Eksporter stemmerne fra en musescore-fil som MIDI og indsæt med passende VST.

\subsubsection{Arrangementskrivning som rytmegruppe/blæser par}
Det kan være en stor opgave at producere et helt arrangement til en revysang, og der skal mange forkellige kompetencer til. Hvis man parrer en person fra rytmegruppen med en blæser er alle kompetencer typisk dækket. På den måde kan man producere arrangementer som par, hvor den ene tager sig af rytmegruppen og den anden tager sig af blæserarrangementet.

\subsubsection{Om korsang i popmusik (please ikke diatoniske parallelle tertser)}
Mange sangere har en impuls til, hvis de bliver bedt om at synge en andenstemme, at synge melodien en terts over. Jeg vil appelere til at man tøjler denne impuls. I popmusik er melodien \textit{meget} vigtig. En stemme der bevæger sig nøjagtigt som melodien, men lige over, vil ofte blive kognitivt registreret \textit{som} melodien, fordi den ligger højere, men med i øvrigt identiske parametre. I alle tilfælde vil den få melodien til at fremstå svagere, hvilket er det modsatte af hvad man gerne vil.

Begrebet \textbf{backing vocals} er mere passende til det, man vil have sine ikke-lead sangere til at levere, end "kor", når det kommer til pop.

Dua Lipas yndlings backing-vocal stemme er \textit{unison} nært efterfulgt af \textit{oktaver}, med tonegentagelser over akkorder på en fjern 3. plads.

Har man brug for inspiration eller eksempler vil jeg anbefale én enkelt sang over alle andre: Ariana Grandes Into You. Den har nok idéer til fede pop-backing vocals til at fylde en hel revy.

\subsubsection{Tekniker til at arrangere backing vocals}
\paragraph{Unison/oktaver}
Backing vocals der er unison med melodien er ikke bare et populært element i popmusik, det er nærmest et definerende genretræk, og burde bruges derefter.
En kvalitet unisone backing vocals er født med, er at leadsangeren kan frasere frit uden melodien forsvinder.
Man skal dog ikke tro at unisone backing vocals er \textit{nemt}, bare fordi det er konceptuelt simpelt.
At styrke melodien uden at tage fokus fra leadsangeren kræver både justering af frasering og artikulation, samt tilpasning til den enkelte leadsanger.
Derudover stiller unison sang høje krav til intonation fra både lead og backup sangere.

Disse betragtninger gør sig også gældende for oktaver.

\paragraph{Rytmiske roller}

\paragraph{Emphasis}

\paragraph{Pentatonic pararllel movements}

\subsubsection{Did you overcook it?}
Nogle gange skal der \textit{cookes}. I sådanne tider må man huske at spørge sig selv:

\textit{Did I overcook it?}

\subsection{Elsk din besætning, kend dit band, udvikl jeres stil}